\section{Conclusion}

This study has investigated the integration of heterogeneous multimodal features with graph-based collaborative filtering architectures to solve the challenges of interaction sparsity and cold-start in personalized fashion recommendation. Through a systematic evaluation of 24 model configurations—spanning three graph neural network (GNN) architectures (BM3, CombiGCN adapted, and FREEDOM), two vision encoders (MobileNetV2 and CLIP), two textual encoders (TF-IDF and BERT), and four similarity and fusion strategies—several key empirical and theoretical conclusions can be drawn:

First, the integration of multimodal attributes with graph convolutions significantly improves personalized recommendation performance over pure collaborative filtering baselines. By propagating aesthetic compatibility and textual semantics along item-item similarity links, graph networks effectively mitigate the cold-start bottleneck for items with sparse interaction logs.

Second, the visual feature extractor plays a critical role in style compatibility representation. Empirical results demonstrate that MobileNetV2 (MBNv2) consistently outperforms or matches the performance of the CLIP vision-language foundation model. While CLIP is optimized for high-level semantic cross-modal alignment (e.g., classifying broad categories like a ``red jacket''), fashion styling compatibility is governed by fine-grained visual details such as local textures, stitching, fabric patterns, and collar lines. MobileNetV2, utilizing local spatial convolutions, preserves these low-level spatial styling features, making it a more suitable encoder for visual garment compatibility.

Third, the selection of a multimodal fusion strategy must be balanced against dataset scale. Late fusion (\texttt{multimodal}) via element-wise average pooling represents a robust, parameter-efficient baseline that consistently yields peak performance. Conversely, trainable weighted attention gating (\texttt{multimodal\_attention}) introduces excessive model complexity and overfitting on sparse, small-scale datasets (such as VCR), degrading the performance of robust models like BM3 and CombiGCN while only slightly assisting weaker baselines (FREEDOM).

Finally, model performance is highly retrieval-depth dependent. BM3, utilizing self-supervised bootstrap representation alignment, represents the state-of-the-art framework for list sizes $K \ge 5$. However, the adapted CombiGCN dual-graph model exhibits superior accuracy for strict top-1 single-item recommendations. These findings demonstrate that visual-textual GNNs are highly effective for personalized fashion recommendation, and their structural components must be carefully optimized for the target application scenario.
